\chapter{Objetivos}
\section{Objetivo general}

Desarrollar un modelo que estime señales rPPG para que clasifique el estrés psicofisiológico en dos 
categorías a partir de video, evaluando su comportamiento ante condiciones sintéticas y 
controladas de degradación lumínica.
%en la base de datos UBFC-Phys, con el fin de aportar 
%evidencia sobre la viabilidad de estas técnicas en escenarios de iluminación no ideal.

\vspace{0.5cm}
\section{Objetivos específicos}
\begin{itemize}

    \item Seleccionar la técnica para extraer señales rPPG identificadas en la revisión 
    de literatura frente a cambios en las condiciones de iluminación.

    % Identificar técnicas mediante la revisión de literatura 
    % %los principales métodos tradicionales
    % % y de aprendizaje profundo utilizados 
    %  para extraer señales rPPG, así como sus limitaciones 
    %  frente a cambios en las condiciones de iluminación, y seleccionar técnicas mas adecuadas 
    %  para el diseño del modelo propuesto.

    \item Preparar un conjunto de datos a partir de videos
    % de reposo (T1) y estrés (T3) de 
    % la base de datos UBFC-Phys, 
    extrayendo la región facial de interés y generando tres 
    condiciones de iluminación.
    % : la original y dos niveles de atenuación sintética, 
    % correspondientes al 50\% y 20\% de luz residual.

    \item Construir el modelo de estimación de rPPG y clasificación del estrés psicofisiológico 
    % (T1 vs. T3),
     integrando los métodos y técnicas seleccionadas en la revisión de literatura y 
    entrenándolo con el conjunto de datos preparado.

    \item Evaluar el desempeño del modelo mediante métricas de error para la estimación de la 
    señal rPPG y métricas de precisión para la clasificación del estrés.
    % , comparando los 
    % resultados obtenidos en las condiciones de iluminación reducida con los de la condición 
    % original y con los valores de referencia reportados en la literatura para UBFC-Phys.

\end{itemize}